We study the original executions of deterministic multitape machines that
transpose a row-major matrix within linear time. An offline replay records
selected run entries and incoming work symbols without fixing a fine head
trace. For fixed resources, uniformly over every positive entry width, all
but exponentially few ambient inputs on which the machine is timely and
correct require two work heads to make linearly many macroscopic returns.
Causal-fragment replay gives original input/output cut expansion and
expansion to off-carrier return predecessors. A guarded projection from an
existing native simulator also gives expansion from every large output set
to original crossings, and capacity certificates on original entry relations,
at canonical widths. These results constrain the actual computation without
assuming that its original trace is shared by many inputs.

Further bounds measure incoming values that escape fixed local
reconstructions, and changes to stored words relative to locally readable
references. A dynamic-carrier argument forces two simultaneous linear-sized
stores. Redundant archives and a cheap historical-feedback control show why
large stores or extensive reuse alone do not prove a time lower bound.
The replay, cut, getter, loader, block-program and actual simulator-transport
arguments are Lean-verified. The dynamic-carrier, archive and physical
monitor constructions have complete written proofs; their remaining Lean
obligations and the consequences conditional on them are explicit.
The question of unrestricted linear-time transposition remains open.
